\chapter*{Conclusion}
\addcontentsline{toc}{chapter}{Conclusion} 

This thesis examines two complementary approaches to automatic machine translation evaluation. The first part focuses on targeted paraphrasing of reference translations as a method to improve the reliability of traditional string-based evaluation metrics. The second part examines the behavior and limitations of modern sentence embedding models with respect to machine translation evaluation. Together, these two approaches highlight the difficulty of measuring semantic equivalence in machine translation.

In addition to the experimental work, the thesis provides a structured overview of machine translation evaluation methods, with a particular focus on human evaluation in the WMT shared tasks. We trace the development from early adequacy and fluency judgments through relative ranking to Direct Assessment, MQM, and more recent approaches such as ESA. This overview clarifies what current evaluation protocols actually measure, as well as their practical limitations, including subjectivity, cost, and limited reproducibility. It also provides an essential context for interpreting automatic metrics, which are typically trained to approximate these human judgments.

Building upon this theoretical foundation, the main contributions of this thesis can be summarized as follows. First, we demonstrate that the reliability of string-based metrics such as \bleu can be systematically improved --- specifically in the context of a morphologically rich language like Czech --- by creating artificial reference sentences closer to machine translation outputs. We show that for Czech, simple substitution-based paraphrasing methods consistently outperform more complex approaches based on statistical or rule-based machine translation systems. This finding suggests that controlled lexical variation is more effective than generating complete paraphrases that differ more significantly in structure. At the same time, we show that large, automatically created paraphrase resources can reduce evaluation quality by introducing noise. Filtering and moving paraphrase processing from evaluation to preprocessing improves results. Finally, we establish that grammatical correctness plays a crucial role: linguistically well-formed paraphrases consistently improve evaluation reliability.

Second, we provide a detailed analysis of sentence embeddings as a semantic evaluation tool. We analyze the geometric properties of both English and multilingual sentence embedding spaces and show that certain models exhibit a clear and interpretable geometric structure, in which simple linguistic transformations correspond to consistent directional shifts. This structure is not a trivial byproduct of architecture, but emerges from the training objective, particularly in models trained with semantic or contrastive signals.

To support further examination of embedding properties, we introduce the \costra dataset of complex Czech sentence transformations, addressing a gap in existing evaluation resources. Our experiments show that even strong, state-of-the-art embedding models continue to rely heavily on surface-level similarity and often fail to capture more complex semantic relations. We also find that strong performance on intrinsic evaluation tasks does not reliably predict performance in machine translation evaluation, despite both being designed to assess semantic similarity.

Taken together, the results lead to a consistent conclusion: a more complex model does not automatically result in better evaluation performance. Simple string-based metrics, when supported by carefully constructed and grammatically sound paraphrases, can perform as well as or better than more complex approaches. Similarly, larger and more general sentence embedding models do not necessarily outperform smaller or more specialized models when applied to machine translation evaluation. These findings demonstrate that evaluation quality depends primarily on how well a method is matched to the specific task, rather than on model complexity alone.

At the same time, several open questions remain. In particular, it is unclear why intrinsic and extrinsic evaluation lead to contradictory conclusions about embedding quality. A deeper understanding of this mismatch would require a more precise characterization of what aspects of meaning are captured by human evaluation and how these relate to current intrinsic benchmarks. We leave this as future work: designing intrinsic evaluation methods that better reflect the requirements of downstream tasks such as machine translation evaluation.

Another promising direction for future work is the deeper exploration of sentence representations with interpretable geometric properties, as observed in our analysis of embedding spaces. Understanding what training conditions lead to these structured spaces---and how to mitigate their failures in areas like logical negation---may contribute to better generalization and more interpretable models.